\chapter{LLM Gateway, Feature Flags \& Self-Healing Documentation}
\label{chap:llm_gateway_featureflags}

% ============================================================
% AUTOR-NOTIZ STATT LERNZIELE
% ============================================================
\autornotiz{
2024: Unser Support-Bot lief auf einem Premium-Modell. Black Friday: Provider Rate-Limit 429, Latenz 15s, Kosten \$4.000/Tag. Kein Fallback, kein Cache, kein Circuit Breaker. Fix in 4h: LiteLLM Gateway + Redis Circuit Breaker + Fallback zu einem Reasoning-Modell + Self-hosted 70B-Modell. Kosten runter auf \$800/Tag, Availability 99.9~\%.
Lektion: Single-Provider ist Fahrlässigkeit. LLM Gateway ist keine Option — es ist Infrastruktur wie Load Balancer oder DB.
}

% ============================================================
% LERNZIELE
% ============================================================
\begin{lernzieleEnv}
\begin{enumerate}
    \item Einen Multi-Provider LLM Gateway mit Unified API, Fallback-Ketten, verteiltem Circuit Breaker und Cost-Aware Routing in Produktion betreiben.
    \item AI Feature Flags für Prompt-Versioning, Model-Switching und Canary-Rollouts mit automatischen Quality Gates (Eval-Score Thresholds) und Instant-Auto-Rollback implementieren.
    \item Self-Healing Documentation Pipelines bauen, die aus Production Logs signifikante Antwort-Drift erkennen, LLM-basiert Dokumentation aktualisieren und via PR + Auto-Merge deployen.
    \item Auto-Eval Dataset Builder aus Production Logs aufsetzen: Clustering (HDBSCAN) → Representative Sampling → LLM-as-Judge Labeling → Golden Dataset Export.
    \item Resiliente LLM-Architektur designen, die Vendor Lock-in vermeidet, Kosten explodieren verhindert und probabilistische Failures graceful degradiert.
\end{enumerate}
\end{lernzieleEnv}

% ============================================================
% MOTIVATION
% ============================================================
\section{Motivation}

Eine LLM-Anwendung in Produktion zu bringen unterscheidet sich grundlegend vom klassischen Software-Deployment. Kapitel 10 (\ref{chap:deployment_architektur}) hat die Basis gelegt: Deployment-Modelle, Latenz, Cost-Performance-Trade-offs, Circuit Breaker. Dieses Kapitel geht den nächsten Schritt — wir bauen die **Produktions-Infrastruktur**, die Vendor Lock-in verhindert, probabilistische Änderungen ohne Deploy ermöglicht und Dokumentation automatisch aktuell hält.

Drei Pain Points treffen jedes LLM-Prod-System:

\begin{enumerate}[leftmargin=*]
    \item \textbf{Vendor Lock-in \& Single Point of Failure} — OpenAI down = dein Produkt down. Keine Verhandlungsmacht bei Preiserhöhungen.
    \item \textbf{Prompt/Model Changes = Code Deploy} — Jede Prompt-Änderung braucht CI/CD, Review, Rollback. Kein Kill-Switch, kein Canary, kein A/B Test.
    \item \textbf{Docs veralten schneller als Code} — Support-Bot antwortet neu, FAQ ist 3 Monate alt. Kunde vertraut Bot nicht. Manuelle Doc-Updates skalieren nicht.
\end{enumerate}

\praxishinweis{LLM Gateway ist keine Option — es ist Infrastruktur wie Load Balancer oder DB. Wer ohne Gateway deployed, zahlt Lehrgeld.}

% ============================================================
% GRUNDLAGEN & ARCHITEKTUR
% ============================================================
\section{Grundlagen \& Architektur}

\subsection{LLM Gateway — Die Abstraktionsschicht über alle Provider}

\paragraph{Was ein Gateway leistet}
Ein LLM Gateway ist die zentrale Abstraktionsschicht zwischen deiner Applikation und den Modell-Providern:

\begin{itemize}
    \item \textbf{Unified API:} Ein Interface (OpenAI-kompatibel) für 100+ Provider (OpenAI, Anthropic, Vertex, Bedrock, vLLM, TGI, Ollama).
    \item \textbf{Request/Response Normalisierung:} Parameter-Mapping, Streaming-Uniformität, Token-Counting über Provider hinweg.
    \item \textbf{Resilience Patterns:} Verteilter Circuit Breaker (Redis-backed), Fallback-Ketten, Rate Limiting (Token/Request/Cost), Request Hedging (Absicherung gegen Tail-Latency).
    \item \textbf{Observability Hooks:} Kosten, Latenz, Token-Usage pro Call — strukturiert geloggt für Prometheus/Langfuse.
\end{itemize}

\paragraph{Architektur-Pattern: Gateway als Sidecar vs. Library}
\begin{itemize}
    \item \textbf{Sidecar (Portkey, Helicone):} Network Hop, aber sprachagnostisch, zentrale Config, Team-unabhängig.
    \item \textbf{Library (LiteLLM Python SDK):} In-Process, Zero Latency, volle Control, Python-only.
    \item \textbf{Empfehlung:} LiteLLM als Library für Python-Services; Sidecar für Polyglot-Stacks.
\end{itemize}

Siehe auch: Deployment Patterns (Sidecar vs. Library) — Kapitel 10, \ref{chap:deployment_architektur}

\realitycheck{Ein Gateway ohne Observability ist blind. Jeder Request muss Trace-ID, Kosten, Latenz, Provider, Fallback-Status, Circuit-Breaker-State liefern. Ohne diese Daten fliegst du blind in Produktion.}

\subsection{AI Feature Flags — Prompt/Model Versioning ohne Deploy}

\paragraph{Warum klassische Feature Flags nicht reichen}
\begin{itemize}
    \item Flags steuern Code-Pfade; AI Flags steuern \textbf{probabilistisches Verhalten} (Prompt, Model, Temperature, Retrieval Config).
    \item Brauchen \textbf{Quality Gates}: "Rollout nur wenn Eval-Score $\ge$ 90~\% Baseline".
    \item Brauchen \textbf{Auto-Rollback}: Metrik bricht ein (Error Rate, Cost, Latency, Refusal) → Flag instantly auf 0~\%.
\end{itemize}

\paragraph{Flag-Typen für AI}
\begin{description}
    \item[Prompt Version Flag] Steuerung welches Prompt-Template (v1 vs v2) geladen wird.
    \item[Model Version Flag] Traffic-Split zwischen Modellen (z.B. 90\% GPT-4o, 10\% Sonnet 4).
    \item[Routing Rule Flag] Aktivierung neuer Routing-Logik (Semantic Router vs. Keyword Router).
    \item[Feature Toggle Flag] Ein/Aus für Features (Semantic Cache, Guardrail, Structured Output).
\end{description}

Siehe auch: Feature Flags in Deployment — Kapitel 10, \ref{chap:deployment_architektur}; Canary Deployments — Kapitel 10

\subsection{Self-Healing Documentation — Docs die sich selbst pflegen}

\paragraph{Problem} Production LLM-Antworten driften von Dokumentation ab. Nutzer vertrauen Docs nicht mehr → Support-Tickets steigen.

\paragraph{Lösung: Pipeline}
\begin{enumerate}
    \item \textbf{Log Ingestion:} Täglich Production Traces (Langfuse/Helicone) holen.
    \item \textbf{Diff Detection:} Embedding-Similarity (text-embedding-3-large) zwischen Prod-Antworten und Doc-Sections. Threshold < 0.85 = Drift.
    \item \textbf{LLM Doc Update:} GPT-4o/Claude-3.5 generiert aktualisiertes Markdown aus QA-Pairs + altem Doc.
    \item \textbf{Validation:} Vale (Style), Markdownlint, LLM-as-Judge (Factuality).
    \item \textbf{PR + Auto-Merge:} GitHub Action erstellt PR, reviewed von Docs-Team, auto-merge bei grünen Checks.
\end{enumerate}

\praxishinweis{Threshold 0.85 Cosine Similarity hat sich bewährt. Darunter = signifikante inhaltliche Änderung. Darunter löst die Pipeline einen Doc-Update aus.}

\subsection{Auto-Eval Dataset Builder — Golden Datasets aus Produktion}

\paragraph{Warum nicht manuell?} Golden Datasets sind teuer (Experten-Labeling), veralten schnell, decken Edge-Cases nicht ab.

\paragraph{Pipeline}
\begin{enumerate}
    \item \textbf{Log Ingestion:} Production Requests/Responses (24h Window).
    \item \textbf{Clustering:} HDBSCAN auf Embeddings (text-embedding-3-large / gte-large-en-v1.5) — findet variable Dichte, markiert Noise.
    \item \textbf{Representative Sampling:} Medoid pro Cluster → diverser, repräsentativer Sample.
    \item \textbf{LLM-as-Judge Labeling:} GPT-4o/Claude-3.5 bewertet (Faithfulness, Relevance, Correctness) mit n=3 + Confidence Scoring.
    \item \textbf{Quality Filter:} Judge Confidence > 0.8 + Consistency Check (3/3 Agreement).
    \item \textbf{Export:} JSONL (OpenAI Eval Format) / HuggingFace Dataset → CI/CD Integration.
\end{enumerate}

Siehe auch: Evaluation Metrics — Kapitel 9, \ref{chap:evaluation}; Observability — Kapitel 15, \ref{chap:mlops_modelllebenszyklus}

% ============================================================
% THEORY DEEP DIVES
% ============================================================
\section{Theory \& Deep Dives}

\subsection{Routing Logic — Intelligence im Gateway}

\paragraph{Routing-Strategien (von einfach bis smart)}
\begin{enumerate}
    \item \textbf{Static Priority:} Feste Fallback-Kette (GPT-4o → Sonnet 4 → Llama-3.3-70b).
    \item \textbf{Intent-Based (Keyword/Classifier):} "Code" → Code-Modell, "Creative" → Creative-Modell.
    \item \textbf{Semantic Routing (Embedding Classifier):} MiniLM/BERT Embedding → Intent Classifier (<10ms, kein LLM Call). Siehe Ch12.
    \item \textbf{Cost-Aware Routing:} \texttt{estimated\_cost < budget\_remaining} → günstiges Modell; Else Premium.
    \item \textbf{Quality-Aware Routing:} Historische Eval-Scores pro Model/Task → Route zum besten Price/Quality Ratio.
    \item \textbf{Weighted Multi-Provider:} 70\% Cheap Model, 30\% Premium für Quality-Sampling.
\end{enumerate}

\paragraph{Request Hedging (Tail Latency Optimization)}
\begin{itemize}
    \item Nach p99-Latency-Threshold (z.B. 3s) Duplicate Request an Backup-Provider.
    \item Erste Response gewinnt; andere wird gecancelt.
    \item Kosten: 2x API Call für Tail Requests (~1-5\% Traffic). Lohnt sich bei p99 > 10s.
\end{itemize}

Siehe auch: Semantic Caching \& Routing — Kapitel 12, \ref{chap:caching}

\subsection{Rollout Strategies — Progressive Delivery für AI}

\paragraph{Canary Rollout mit Quality Gates}
\begin{enumerate}
    \item \textbf{Pre-Flight Eval:} Candidate Config (Prompt v2, Model X) vs. Golden Dataset → Score $\ge$ Baseline × Threshold.
    \item \textbf{Shadow Traffic:} 100~\% Traffic an Current + Candidate parallel, nur Metriken vergleichen (kein User Impact).
    \item \textbf{Progressive Rollout:} 5~\% → 25~\% → 50~\% → 100~\% (konfigurierbare Steps + Intervalle).
    \item \textbf{Continuous Monitoring:} Auto-Rollback bei Metrik-Violation (Error Rate > 5~\%, P99 Latency > 3× Baseline, Cost > 1.5× Baseline, Refusal Rate > 10~\%).
    \item \textbf{Full Rollout:} Flag auf 100~\%, Baseline aktualisieren.
\end{enumerate}

\paragraph{Argo Rollouts / Flagger Integration}
\begin{itemize}
    \item Kubernetes-native Progressive Delivery.
    \item AnalysisTemplate definiert Metriken + Thresholds.
    \item Automatischer Rollback/Abort bei Violation.
\end{itemize}

\praxishinweis{Shadow Traffic (Schattenverkehr) ist sicherer als direkter Canary. 100~\% Requests an Current + Candidate parallel. User sieht nur Current Response. Candidate Metriken werden gesammelt ohne User Impact. Erst nach 24h Shadow-Analyse (keine Regressionen) startet Canary Rollout.}

\subsection{Clustering Algorithms für Auto-Eval — HDBSCAN Deep Dive}

\paragraph{Warum HDBSCAN nicht K-Means?}
\begin{itemize}
    \item K-Means: Fixes k, kugelförmige Clusters, sensitiv zu Outliers.
    \item HDBSCAN: Hierarchisch, variable Dichte, \textbf{Noise Points explizit}, kein k nötig.
    \item Production Logs haben variable Cluster-Dichten (FAQs = dicht, Edge Cases = sparse) + viel Rauschen.
\end{itemize}

\paragraph{Parameter Tuning für LLM Logs}
\begin{itemize}
    \item \texttt{min\_cluster\_size}: 10–20 (kleinere Clusters = granularere Topics).
    \item \texttt{min\_samples}: 1–5 (niedriger = mehr Clusters, höher = konservativer).
    \item \texttt{metric}: 'cosine' für Embeddings.
    \item \texttt{cluster\_selection\_method}: 'eom' (Excess of Mass) oder 'leaf' (feingranularer).
\end{itemize}

\paragraph{Representative Sampling}
\begin{itemize}
    \item \textbf{Medoid:} Punkt mit minimaler durchschnittlicher Distanz zu allen Cluster-Mitgliedern (robust).
    \item \textbf{Centroid:} Mittelwert der Embeddings (schneller, aber sensitiv zu Outliers).
    \item \textbf{Max-Min Diversity:} Iterativ am weitesten entfernten Punkt wählen (für Coverage).
\end{itemize}

% ============================================================
% PRAXIS & IMPLEMENTATION (CODE-ZENTRIERT)
% ============================================================
\section{Praxis \& Implementation}

\subsection{Gateway Core — Multi-Provider Router mit Resilience}

\begin{lstlisting}[language=Python, caption={LLM Gateway Core: Multi-Provider Router, Circuit Breaker, Fallback Chains, Rate Limiting}, label=lst:ch17_gateway_core]
from __future__ import annotations
import asyncio
import time
from dataclasses import dataclass, field
from enum import Enum
from typing import Optional
import redis.asyncio as redis
import litellm
from litellm import acompletion

class CircuitState(Enum):
    CLOSED = "closed"
    OPEN = "open"
    HALF_OPEN = "half_open"

@dataclass
class ModelConfig:
    name: str
    provider: str
    tags: list[str] = field(default_factory=list)  # fast, cheap, reasoning, self-hosted
    cost_per_1k_input: float = 0.0
    cost_per_1k_output: float = 0.0
    max_tokens: int = 4096

@dataclass
class CircuitBreakerState:
    state: CircuitState = CircuitState.CLOSED
    failures: int = 0
    last_failure: float = 0.0
    successes: int = 0

class DistributedCircuitBreaker:
    """Redis-backed Circuit Breaker für verteilte Instanzen."""
    def __init__(self, redis_client: redis.Redis, model_name: str,
                 failure_threshold: int = 5, recovery_timeout: float = 30.0,
                 half_open_max_calls: int = 3):
        self.redis = redis_client
        self.key = f"circuit_breaker:{model_name}"
        self.failure_threshold = failure_threshold
        self.recovery_timeout = recovery_timeout
        self.half_open_max_calls = half_open_max_calls

    async def _load(self) -> CircuitBreakerState:
        data = await self.redis.hgetall(self.key)
        if not data:
            return CircuitBreakerState()
        return CircuitBreakerState(
            state=CircuitState(data.get(b"state", b"closed").decode()),
            failures=int(data.get(b"failures", 0)),
            last_failure=float(data.get(b"last_failure", 0)),
            successes=int(data.get(b"successes", 0)),
        )

    async def _save(self, state: CircuitBreakerState):
        await self.redis.hset(self.key, mapping={
            "state": state.state.value,
            "failures": state.failures,
            "last_failure": state.last_failure,
            "successes": state.successes,
        })
        await self.redis.expire(self.key, int(self.recovery_timeout * 2))

    async def can_execute(self) -> bool:
        state = await self._load()
        if state.state == CircuitState.CLOSED:
            return True
        if state.state == CircuitState.OPEN:
            if time.time() - state.last_failure > self.recovery_timeout:
                state.state = CircuitState.HALF_OPEN
                state.successes = 0
                await self._save(state)
                return True
            return False
        # HALF_OPEN
        return state.successes < self.half_open_max_calls

    async def record_success(self):
        state = await self._load()
        if state.state == CircuitState.HALF_OPEN:
            state.successes += 1
            if state.successes >= self.half_open_max_calls:
                state.state = CircuitState.CLOSED
                state.failures = 0
        await self._save(state)

    async def record_failure(self):
        state = await self._load()
        state.failures += 1
        state.last_failure = time.time()
        if state.failures >= self.failure_threshold:
            state.state = CircuitState.OPEN
        await self._save(state)

class TokenBucketRateLimiter:
    """Redis-backed Token Bucket für User-Level + Global Rate Limiting."""
    def __init__(self, redis_client: redis.Redis,
                 user_rpm: int = 60, user_tpm: int = 100_000,
                 global_rpm: int = 1000):
        self.redis = redis_client
        self.user_rpm = user_rpm
        self.user_tpm = user_tpm
        self.global_rpm = global_rpm

    async def check_limit(self, user_id: str, estimated_tokens: int = 1000) -> bool:
        now = time.time()
        minute_key = int(now // 60)

        # User Requests/Min
        user_req_key = f"ratelimit:req:{user_id}:{minute_key}"
        user_req = await self.redis.incr(user_req_key)
        if user_req == 1:
            await self.redis.expire(user_req_key, 120)
        if user_req > self.user_rpm:
            return False

        # User Tokens/Min (Token Bucket Approximation)
        user_tok_key = f"ratelimit:tok:{user_id}:{minute_key}"
        user_tok = await self.redis.incrby(user_tok_key, estimated_tokens)
        if user_tok == estimated_tokens:
            await self.redis.expire(user_tok_key, 120)
        if user_tok > self.user_tpm:
            return False

        # Global Requests/Min
        global_key = f"ratelimit:req:global:{minute_key}"
        global_req = await self.redis.incr(global_key)
        if global_req == 1:
            await self.redis.expire(global_key, 120)
        if global_req > self.global_rpm:
            return False

        return True

@dataclass
class FallbackChain:
    primary: str
    fallbacks: list[str]  # Model names in priority order

class ModelRegistry:
    def __init__(self):
        self.models: dict[str, ModelConfig] = {}
        self.fallback_chains: dict[str, FallbackChain] = {}

    def register(self, config: ModelConfig):
        self.models[config.name] = config

    def set_fallback_chain(self, primary: str, fallbacks: list[str]):
        self.fallback_chains[primary] = FallbackChain(primary=primary, fallbacks=fallbacks)

    def get_fallback_chain(self, model_name: str) -> FallbackChain:
        return self.fallback_chains.get(model_name, FallbackChain(primary=model_name, fallbacks=[]))

class LLMGateway:
    """Produktionsreifer Multi-Provider Gateway."""
    def __init__(self, redis_url: str = "redis://localhost:6379"):
        self.redis = redis.from_url(redis_url, decode_responses=True)
        self.registry = ModelRegistry()
        self.circuit_breakers: dict[str, DistributedCircuitBreaker] = {}
        self.rate_limiter = TokenBucketRateLimiter(self.redis)

        # Default Models registrieren
        self._setup_default_models()

    def _setup_default_models(self):
        self.registry.register(ModelConfig(
            name="premium", provider="openai", model="gpt-4o",
            tags=["premium", "reasoning"], cost_per_1k_input=2.50, cost_per_1k_output=10.00
        ))
        self.registry.register(ModelConfig(
            name="fast_cheap", provider="openai", model="gpt-4o-mini",
            tags=["fast", "cheap"], cost_per_1k_input=0.15, cost_per_1k_output=0.60
        ))
        self.registry.register(ModelConfig(
            name="reasoning", provider="anthropic", model="claude-3-5-sonnet",
            tags=["premium", "reasoning"], cost_per_1k_input=3.00, cost_per_1k_output=15.00
        ))
        self.registry.register(ModelConfig(
            name="self_hosted_70b", provider="vllm", model="meta-llama/Llama-3.3-70B-Instruct",
            tags=["self-hosted", "cheap", "private"], cost_per_1k_input=0.0, cost_per_1k_output=0.0
        ))

        # Fallback Chains
        self.registry.set_fallback_chain("premium", ["reasoning", "fast_cheap", "self_hosted_70b"])
        self.registry.set_fallback_chain("fast_cheap", ["self_hosted_70b", "reasoning"])
        self.registry.set_fallback_chain("reasoning", ["premium", "fast_cheap"])
        self.registry.set_fallback_chain("self_hosted_70b", ["fast_cheap", "premium"])

        # Circuit Breaker für jedes Modell
        for name in self.registry.models:
            self.circuit_breakers[name] = DistributedCircuitBreaker(self.redis, name)

    async def chat_completion(self, user_id: str, messages: list[dict],
                              model: str = "auto", temperature: float = 0.1,
                              max_tokens: int = 1024) -> dict:
        """Haupteinstieg: Routing, CB, Fallback, Rate Limit, Cost Tracking."""
        # 1. Model Resolution
        target_model = self._resolve_model(model, messages)
        config = self.registry.models[target_model]

        # 2. Rate Limit Check
        if not await self.rate_limiter.check_limit(user_id, max_tokens):
            return {"error": "rate_limit_exceeded", "retry_after_seconds": 60}

        # 3. Circuit Breaker Check
        cb = self.circuit_breakers[target_model]
        if not await cb.can_execute():
            # Try Fallback
            chain = self.registry.get_fallback_chain(target_model)
            for fallback in chain.fallbacks:
                fb_cb = self.circuit_breakers[fallback]
                if await fb_cb.can_execute():
                    target_model = fallback
                    config = self.registry.models[target_model]
                    cb = fb_cb
                    break
            else:
                return {"error": "all_models_unavailable", "circuit_breakers_open": True}

        # 4. Execute with Retry
        for attempt in range(3):
            try:
                start = time.time()
                response = await acompletion(
                    model=config.name,
                    messages=messages,
                    temperature=temperature,
                    max_tokens=max_tokens,
                )
                latency_ms = (time.time() - start) * 1000

                # Cost Calculation
                usage = response.usage
                cost = (usage.prompt_tokens / 1000 * config.cost_per_1k_input +
                        usage.completion_tokens / 1000 * config.cost_per_1k_output)

                await cb.record_success()

                return {
                    "response": response.choices[0].message.content,
                    "model": target_model,
                    "provider": config.provider,
                    "tokens": {"input": usage.prompt_tokens, "output": usage.completion_tokens},
                    "cost_usd": cost,
                    "latency_ms": latency_ms,
                    "fallback_used": target_model != model,
                }

            except Exception as e:
                await cb.record_failure()
                if attempt == 2:
                    return {"error": "max_retries_exceeded", "details": str(e)}
                await asyncio.sleep(2 ** attempt)  # Exponential backoff

        return {"error": "unexpected_failure"}

    def _resolve_model(self, model: str, messages: list[dict]) -> str:
        if model != "auto":
            return model
        # Simple heuristic: short context -> fast_cheap, long -> premium
        total_chars = sum(len(m.get("content", "")) for m in messages)
        if total_chars < 2000:
            return "fast_cheap"
        return "premium"
\end{lstlisting}

\realitycheck{In-Memory Circuit Breaker ist für Produktion unbrauchbar. State im Prozess → bei Restart/Scale-out verloren. Bei 3 Replicas hat jeder eigenen State → Inconsistente CB-Entscheidungen. Verteilter CB (Redis) teilt State: Alle Instanzen sehen gleichen State, Recovery koordiniert, keine "Split Brain" Failures.}

\subsection{Feature Flag Engine mit Quality Gates \& Auto-Rollback}

\begin{lstlisting}[language=Python, caption={AI Feature Flag Engine: Unleash Integration, Quality Gate, Auto-Rollback, Rollout Controller}, label=lst:ch17_feature_flags]
from dataclasses import dataclass, field
from typing import Literal
import asyncio
from unleash_client import UnleashClient

@dataclass
class AIFlag:
    """Feature Flag für AI-Systeme mit Quality Gates und Auto-Rollback."""
    name: str
    flag_type: Literal["prompt_version", "model_version", "routing_rule", "feature_toggle"]
    eval_dataset: str          # Pfad zu Golden Dataset (JSONL)
    quality_threshold: float   # z.B. 0.90 = 90% der Baseline
    rollout_steps: list[int] = field(default_factory=lambda: [5, 25, 50, 100])
    rollout_interval_hours: int = 24
    rollback_metrics: dict = field(default_factory=lambda: {
        "error_rate": 0.05,           # > 5% in 5min
        "p99_latency_multiplier": 3.0, # > 3x Baseline in 10min
        "cost_multiplier": 1.5,        # > 1.5x Baseline in 60min
        "refusal_rate": 0.10,          # > 10% in 10min
    })

class QualityGateEvaluator:
    """Pre-Flight Evaluation gegen Golden Dataset."""

    def __init__(self, evaluator_model: str = "gpt-4o"):
        self.evaluator_model = evaluator_model

    async def evaluate(self, flag: AIFlag, candidate_config: dict) -> tuple[bool, dict]:
        """
        Returns: (passed, metrics_dict)
        """
        from evaluation import run_golden_dataset  # Ch9
        baseline_metrics = await run_golden_dataset(flag.eval_dataset, candidate_config)
        # Simplified: Compare Faithfulness/Relevance vs Baseline
        faithfulness = baseline_metrics.get("faithfulness", 0.0)
        relevance = baseline_metrics.get("relevance", 0.0)
        quality_score = (faithfulness + relevance) / 2
        passed = quality_score >= flag.quality_threshold
        return passed, {"faithfulness": faithfulness, "relevance": relevance, "score": quality_score}


class AutoRollbackMonitor:
    """Hintergrund-Task: Prüft Prod-Metriken, triggert Rollback bei Violation."""

    def __init__(self, unleash: UnleashClient, prometheus_url: str):
        self.unleash = unleash
        self.prometheus_url = prometheus_url

    async def check_and_rollback(self, flag: AIFlag) -> bool:
        # Query Prometheus for each metric
        violations = []
        for metric, threshold in flag.rollback_metrics.items():
            current = await self._query_metric(metric)
            baseline = await self._get_baseline(metric)
            if self._is_violation(metric, current, baseline, threshold):
                violations.append(f"{metric}: {current:.2f} > {threshold}x baseline ({baseline:.2f})")

        if violations:
            # Kill Switch: Flag auf 0%
            await self.unleash.update_flag(flag.name, {"percentage": 0})
            await self._alert(violations)
            return True
        return False

    async def _query_metric(self, metric: str) -> float:
        # PromQL query per metric
        pass

    async def _get_baseline(self, metric: str) -> float:
        # Load baseline from config
        pass

    def _is_violation(self, metric: str, current: float, baseline: float, threshold) -> bool:
        if isinstance(threshold, float):
            return current > baseline * threshold
        return False

    async def _alert(self, violations: list[str]):
        # Slack/PagerDuty
        pass


class RolloutController:
    """Orchestriert: Quality Gate → Progressive Steps → Monitoring → Completion."""

    def __init__(self, unleash: UnleashClient, quality_gate: QualityGateEvaluator,
                 rollback_monitor: AutoRollbackMonitor):
        self.unleash = unleash
        self.quality_gate = quality_gate
        self.rollback_monitor = rollback_monitor

    async def rollout(self, flag: AIFlag, candidate_config: dict):
        # 1. Pre-Flight Quality Gate
        passed, metrics = await self.quality_gate.evaluate(flag, candidate_config)
        if not passed:
            raise ValueError(f"Quality Gate failed: {metrics}")

        # 2. Progressive Rollout
        for step_pct in flag.rollout_steps:
            await self.unleash.update_flag(flag.name, {"percentage": step_pct})
            print(f"Rollout {flag.name}: {step_pct}%")

            # Monitor für rollout_interval_hours
            end_time = asyncio.get_event_loop().time() + flag.rollout_interval_hours * 3600
            while asyncio.get_event_loop().time() < end_time:
                await asyncio.sleep(60)
                if await self.rollback_monitor.check_and_rollback(flag):
                    raise RuntimeError(f"Auto-Rollback triggered for {flag.name}")

        # 3. Full Rollout
        await self.unleash.update_flag(flag.name, {"percentage": 100})
        print(f"Rollout {flag.name} completed: 100%")
\end{lstlisting}

\paragraph{Key Classes}
\begin{itemize}
    \item \texttt{AIFlag} — Dataclass mit Flag-Type, Eval-Dataset, Threshold, Rollout-Steps, Rollback-Metriken.
    \item \texttt{QualityGateEvaluator} — Pre-Flight Eval gegen Golden Dataset (relative Thresholds).
    \item \texttt{AutoRollbackMonitor} — Hintergrund-Task, prüft Prod-Metriken (Prometheus/Langfuse) jede Minute.
    \item \texttt{RolloutController} — Orchestriert: Quality Gate → Progressive Steps → Monitoring → Completion.
\end{itemize}

\praxishinweis{Quality Gate Threshold = 0.85–0.95 je Kritikalität. Support-Bot: 0.90. Premium-QA: 0.95. Nie 1.0 — das blockiert jeden Rollout. Relative Thresholds (Baseline × 0.9) sind robuster als absolute, da sich Baseline mit Modell-Updates verschiebt.}

\subsection{Self-Healing Documentation Pipeline (GitHub Actions)}

\begin{lstlisting}[language=bash, caption={GitHub Actions Workflow: Detect Changes → Generate Docs → Validate → PR → Auto-Merge}, label=lst:ch17_self_healing_docs]
name: Self-Healing Documentation
on:
  schedule:
    - cron: '0 3 * * *'  # Daily 03:00 UTC
  workflow_dispatch:

jobs:
  detect-changes:
    runs-on: ubuntu-latest
    outputs:
      drift_detected: ${{ steps.detect.outputs.drift_detected }}
      changed_sections: ${{ steps.detect.outputs.changed_sections }}
    steps:
      - uses: actions/checkout@v4
      - name: Set up Python
        uses: actions/setup-python@v5
        with: { python-version: '3.11' }
      - name: Install deps
        run: pip install langfuse openai hdbscan scikit-learn numpy
      - name: Run drift detection
        id: detect
        env:
          LANGFUSE_PUBLIC_KEY: ${{ secrets.LANGFUSE_PUBLIC_KEY }}
          LANGFUSE_SECRET_KEY: ${{ secrets.LANGFUSE_SECRET_KEY }}
          OPENAI_API_KEY: ${{ secrets.OPENAI_API_KEY }}
        run: python scripts/detect_doc_changes.py

  generate-docs:
    needs: detect-changes
    if: needs.detect-changes.outputs.drift_detected == 'true'
    runs-on: ubuntu-latest
    steps:
      - uses: actions/checkout@v4
      - name: Set up Python
        uses: actions/setup-python@v5
        with: { python-version: '3.11' }
      - name: Generate doc updates
        id: generate
        env:
          OPENAI_API_KEY: ${{ secrets.OPENAI_API_KEY }}
        run: python scripts/generate_doc_update.py
      - name: Upload generated docs
        uses: actions/upload-artifact@v4
        with: { name: updated-docs, path: docs/_generated/ }

  validate:
    needs: generate-docs
    runs-on: ubuntu-latest
    steps:
      - uses: actions/checkout@v4
      - name: Download generated docs
        uses: actions/download-artifact@v4
        with: { name: updated-docs, path: docs/_generated/ }
      - name: Run validation
        run: |
          vale docs/_generated/
          markdownlint docs/_generated/
          python scripts/factuality_check.py

  create-pr:
    needs: validate
    runs-on: ubuntu-latest
    steps:
      - uses: actions/checkout@v4
      - name: Download generated docs
        uses: actions/download-artifact@v4
        with: { name: updated-docs, path: docs/_generated/ }
      - name: Create PR
        uses: peter-evans/create-pull-request@v6
        with:
          branch: self-healing-docs-${{ github.run_id }}
          title: "docs: auto-update from production drift"
          body: |
            Automatisch generierte Dokumentations-Updates basierend auf Produktions-Drift.
            Geänderte Sektionen: ${{ needs.detect-changes.outputs.changed_sections }}
          labels: documentation, automated
          draft: false

  auto-merge:
    needs: create-pr
    if: github.event.pull_request.labels.*.name contains 'automated'
    runs-on: ubuntu-latest
    steps:
      - name: Enable auto-merge
        uses: peter-evans/enable-pull-request-automerge@v3
        with:
          pull-request-number: ${{ needs.create-pr.outputs.pull-request-number }}
          merge-method: squash
\end{lstlisting}

\begin{lstlisting}[language=Python, caption={Drift Detection: Embedding Similarity zwischen Prod-Antworten und Doc-Sections}, label=lst:ch17_detect_changes]
import numpy as np
from langfuse import Langfuse
from openai import OpenAI
from sklearn.metrics.pairwise import cosine_similarity

langfuse = Langfuse()
client = OpenAI()

def get_production_answers(hours: int = 24) -> list[dict]:
    """Holt Traces aus Langfuse, extrahiert User-Query + Model-Answer."""
    traces = langfuse.trace.list(
        from_timestamp=datetime.now() - timedelta(hours=hours),
        limit=5000,
    )
    qa_pairs = []
    for trace in traces:
        if trace.output and trace.input:
            qa_pairs.append({
                "query": trace.input.get("query", ""),
                "answer": trace.output.get("content", ""),
                "trace_id": trace.id,
            })
    return qa_pairs

def get_document_sections(doc_path: str) -> list[dict]:
    """Parst Markdown-Docs in Sections (Heading + Content)."""
    sections = []
    with open(doc_path) as f:
        content = f.read()
    # Simple parser: split by ## headings
    parts = content.split("\n## ")
    for part in parts[1:]:  # Skip first (before first ##)
        lines = part.split("\n")
        heading = lines[0]
        body = "\n".join(lines[1:]).strip()
        if body:
            sections.append({"heading": heading, "content": body})
    return sections

def detect_drift(qa_pairs: list[dict], doc_sections: list[dict],
                 threshold: float = 0.85) -> list[dict]:
    """Embedding-basierte Drift-Detection."""
    # Embed all QA answers
    answers = [q["answer"] for q in qa_pairs]
    answer_embeddings = embed_batch(answers)

    # Embed all doc sections
    doc_texts = [s["content"] for s in doc_sections]
    doc_embeddings = embed_batch(doc_texts)

    # For each QA pair, find max similarity to any doc section
    drift_results = []
    sim_matrix = cosine_similarity(answer_embeddings, doc_embeddings)
    for i, qa in enumerate(qa_pairs):
        max_sim = sim_matrix[i].max()
        best_section_idx = sim_matrix[i].argmax()
        if max_sim < threshold:
            drift_results.append({
                "trace_id": qa["trace_id"],
                "query": qa["query"],
                "answer": qa["answer"],
                "best_match_section": doc_sections[best_section_idx]["heading"],
                "similarity": max_sim,
            })
    return drift_results

def embed_batch(texts: list[str]) -> np.ndarray:
    """Batch embedding via OpenAI text-embedding-3-large."""
    embeddings = []
    for i in range(0, len(texts), 100):
        batch = texts[i:i+100]
        resp = client.embeddings.create(
            model="text-embedding-3-large",
            input=batch,
        )
        embeddings.extend([np.array(d.embedding) for d in resp.data])
    return np.array(embeddings)

if __name__ == "__main__":
    qa_pairs = get_production_answers(24)
    doc_sections = get_document_sections("docs/support_faq.md")
    drift = detect_drift(qa_pairs, doc_sections, 0.85)

    # Output for next step
    import json
    with open("drift_report.json", "w") as f:
        json.dump(drift, f)
    print(f"Drift detected in {len(drift)} cases")

    # Set GitHub Actions output
    import os
    with open(os.environ["GITHUB_OUTPUT"], "a") as f:
        f.write(f"drift_detected={str(len(drift) > 0).lower()}\n")
        f.write(f"changed_sections={json.dumps([d['best_match_section'] for d in drift])}\n")
\end{lstlisting}

\begin{lstlisting}[language=Python, caption={LLM-basierte Doc-Aktualisierung aus QA-Pairs}, label=lst:ch17_generate_docs]
import json
from openai import OpenAI

client = OpenAI()

def generate_doc_update(drift_report_path: str, doc_path: str) -> str:
    """Generiert aktualisiertes Markdown für gedriftete Sektionen."""
    with open(drift_report_path) as f:
        drift = json.load(f)
    with open(doc_path) as f:
        current_doc = f.read()

    # Group drift by section
    from collections import defaultdict
    by_section = defaultdict(list)
    for d in drift:
        by_section[d["best_match_section"]].append(d)

    updated_sections = {}
    for section, cases in by_section.items():
        # Build prompt: QA pairs + current section
        qa_pairs_text = "\n\n".join(
            f"User: {c['query']}\nBot: {c['answer']}" for c in cases
        )
        prompt = f"""Du aktualisierst Dokumentation basierend auf echten Produktions-QA-Pairs.

AKTUELLE SEKTION:
{get_section_content(current_doc, section)}

NEUE QA-PAIRS AUS PRODUKTION:
{qa_pairs_text}

AUFGABE: Aktualisiere die Sektion so, dass sie die neuen Antworten widerspiegelt.
REGELN:
- Nur Fakten aus QA-Pairs übernehmen, keine Halluzinationen
- Bestehende Struktur/Ton beibehalten
- Ausgabe: NUR das aktualisierte Markdown für diese Sektion
"""
        resp = client.chat.completions.create(
            model="gpt-4o",
            messages=[{"role": "user", "content": prompt}],
            temperature=0.1,
        )
        updated_sections[section] = resp.choices[0].message.content

    # Rebuild full doc
    new_doc = rebuild_document(current_doc, updated_sections)
    return new_doc

def get_section_content(doc: str, heading: str) -> str:
    # Extract section by heading
    pass

def rebuild_document(doc: str, updated: dict) -> str:
    # Replace sections
    pass

if __name__ == "__main__":
    new_doc = generate_doc_update("drift_report.json", "docs/support_faq.md")
    with open("docs/_generated/support_faq.md", "w") as f:
        f.write(new_doc)
\end{lstlisting}

\paragraph{Validierungspipeline}
\begin{enumerate}
    \item \textbf{Style Check:} Vale (Style-Guide), Markdownlint (Syntax).
    \item \textbf{Factuality Check:} LLM-as-Judge prüft: "Stimmt generierter Content mit QA-Pairs überein?"
    \item \textbf{Link Check:} Keine toten Links.
    \item \textbf{Human-in-the-Loop:} PR requires Reviewer Approval (Docs-Team) vor Auto-Merge. Auto-Merge nur bei allen Checks grün.
\end{enumerate}

\warnung{Halluzinationen in Docs sind gefährlich. Mehrstufige Validation ist Pflicht: (1) Source Grounding im Prompt, (2) Vale/Markdownlint, (3) LLM-as-Judge Factuality, (4) Human Review. Auto-Merge nur bei allen grün.}

\subsection{Auto-Eval Dataset Builder}

\begin{lstlisting}[language=Python, caption={Auto-Eval Pipeline: Log Ingestion → HDBSCAN Clustering → Medoid Sampling → LLM-as-Judge Labeling → Quality Filter → Dataset Export}, label=lst:ch17_auto_eval_builder]
import json
import numpy as np
import hdbscan
from langfuse import Langfuse
from openai import OpenAI
from sklearn.metrics.pairwise import cosine_similarity

langfuse = Langfuse()
client = OpenAI()

def fetch_production_logs(days: int = 7) -> list[dict]:
    """Holt alle Traces der letzten N Tage."""
    traces = langfuse.trace.list(
        from_timestamp=datetime.now() - timedelta(days=days),
        limit=20000,
    )
    return [
        {
            "trace_id": t.id,
            "query": t.input.get("query", ""),
            "answer": t.output.get("content", ""),
            "metadata": t.metadata,
        }
        for t in traces if t.output and t.input
    ]

def cluster_queries(logs: list[dict]) -> tuple[np.ndarray, np.ndarray]:
    """HDBSCAN Clustering auf Query-Embeddings."""
    queries = [log["query"] for log in logs]
    embeddings = embed_batch(queries)

    clusterer = hdbscan.HDBSCAN(
        min_cluster_size=15,
        min_samples=3,
        metric='cosine',
        cluster_selection_method='eom',
    )
    labels = clusterer.fit_predict(embeddings)
    return labels, embeddings

def sample_medoids(logs: list[dict], labels: np.ndarray,
                   embeddings: np.ndarray, max_samples: int = 200) -> list[dict]:
    """Wählt Medoid (zentralster Punkt) pro Cluster."""
    from collections import defaultdict
    cluster_indices = defaultdict(list)
    for i, label in enumerate(labels):
        if label >= 0:  # Ignore noise (-1)
            cluster_indices[label].append(i)

    samples = []
    for label, indices in cluster_indices.items():
        cluster_embs = embeddings[indices]
        # Medoid: point with min average distance to others
        dist_matrix = 1 - cosine_similarity(cluster_embs)
        medoid_idx = indices[dist_matrix.mean(axis=1).argmin()]
        samples.append(logs[medoid_idx])

    # Limit total samples
    if len(samples) > max_samples:
        # Diversity sampling
        samples = diversify_samples(samples, max_samples)
    return samples

def label_with_llm(samples: list[dict], n_judges: int = 3) -> list[dict]:
    """LLM-as-Judge Labeling mit Confidence Scoring."""
    judges = ["gpt-4o", "claude-3-5-sonnet", "gemini-1.5-pro"]
    labeled = []

    for sample in samples:
        judge_scores = []
        for judge in judges[:n_judges]:
            prompt = f"""Bewerte die Antwort auf Skala 1-10.
Frage: {sample['query']}
Antwort: {sample['answer']}
Kriterien: Faithfulness, Relevance, Correctness
Gib JSON zurück: {{\"faithfulness\": 1-10, \"relevance\": 1-10, \"correctness\": 1-10, \"confidence\": 0.0-1.0}}"""
            resp = client.chat.completions.create(
                model=judge,
                messages=[{"role": "user", "content": prompt}],
                temperature=0.0,
                response_format={"type": "json_object"},
            )
            judge_scores.append(json.loads(resp.choices[0].message.content))

        # Aggregate
        avg_scores = {k: np.mean([s[k] for s in judge_scores]) for k in judge_scores[0]}
        confidence = avg_scores.pop("confidence")
        consistency = len(set(tuple(sorted(s.items())) for s in judge_scores)) == 1

        if confidence > 0.8 and consistency:
            sample.update({
                "golden_scores": avg_scores,
                "judge_confidence": confidence,
                "judge_consistency": consistency,
            })
            labeled.append(sample)

    return labeled

def export_dataset(labeled: list[dict], version: int) -> str:
    """Export als JSONL im OpenAI Eval Format."""
    output_path = f"eval_datasets/golden_v{version}.jsonl"
    with open(output_path, "w") as f:
        for item in labeled:
            f.write(json.dumps({
                "id": item["trace_id"],
                "input": item["query"],
                "expected": {"reference_answer": item["answer"]},
                "metadata": {
                    "faithfulness": item["golden_scores"]["faithfulness"],
                    "relevance": item["golden_scores"]["relevance"],
                    "correctness": item["golden_scores"]["correctness"],
                    "source": "auto_eval",
                    "version": version,
                }
            }) + "\n")
    return output_path

if __name__ == "__main__":
    logs = fetch_production_logs(7)
    labels, embeddings = cluster_queries(logs)
    samples = sample_medoids(logs, labels, embeddings, 200)
    labeled = label_with_llm(samples)
    path = export_dataset(labeled, version=2)  # Auto-increment
    print(f"Generated {len(labeled)} golden cases → {path}")
\end{lstlisting}

\paragraph{Pipeline-Schritte im Detail}
\begin{enumerate}
    \item \textbf{Log Ingestion:} Langfuse/Helicone SDK → letzte 7 Tage Traces.
    \item \textbf{Clustering:} HDBSCAN auf Embeddings — findet Topics ohne festes k, markiert Noise (seltene/fehlerhafte Queries).
    \item \textbf{Sampling:} Medoid pro Cluster (robust) → 200 Samples gesamt.
    \item \textbf{Labeling:} GPT-4o/Claude-3.5/Gemini, n=3, Confidence Scoring.
    \item \textbf{Quality Filter:} Confidence > 0.8 + Consistency 3/3 Agreement.
    \item \textbf{Export:} JSONL v\{N+1\} versioniert, Baseline aktualisiert.
\end{enumerate}

\realitycheck{Golden Datasets veralten in Wochen. Production Logs sind kostenlos, repräsentativ, decken echte Edge Cases ab. Auto-Eval Pipeline läuft wöchentlich, neues Dataset versioniert (v1, v2, …), Baseline aktualisiert. Ohne Automatisierung evalust du gegen 6 Monate alte Daten.}

% ============================================================
% BEST PRACTICES
% ============================================================
\section{Best Practices}

\begin{enumerate}[leftmargin=*]
    \item \textbf{Gateway First, nicht Last.} Baue den Gateway \textit{vor} dem ersten LLM-Call. Nachträgliches Einbauen kostet 10x mehr.
    \item \textbf{Unified API = OpenAI Format.} Alle Provider auf OpenAI Chat Completions Interface mappen. Vermeidet Vendor-SDK-Lock-in im App-Code.
    \item \textbf{Cost Tracking pro Request.} Jeder Gateway-Response liefert \texttt{cost\_usd}, \texttt{tokens}, \texttt{latency\_ms}. Ohne Cost-Metric fliegst du blind.
    \item \textbf{Feature Flags für alles Probabilistische.} Prompt, Model, Temperature, Retrieval-Config, Guardrails — alles flag-gesteuert.
    \item \textbf{Quality Gate vor Rollout.} Nie ohne Eval-Score deployen. Baseline = Current Production Config. Threshold = 0.85–0.95 je Kritikalität.
    \item \textbf{Auto-Rollback ist Pflicht, nicht Nice-to-Have.} Definiere 4–5 Key Metriken (Error Rate, P99 Latency, Cost/Query, Refusal Rate, Quality Score) mit Multiplikatoren.
\item \textbf{Shadow Traffic vor Canary.} 100~\% Shadow für 24h deckt Regressionen auf, ohne User Impact.
\item \textbf{Self-Healing Docs als CI Job.} Täglich 03:00 UTC. Threshold 0.85 Cosine Similarity. Vale + Markdownlint + LLM-Judge als Quality Gates.
\item \textbf{Golden Dataset Refresh Weekly.} Auto-Eval Pipeline läuft wöchentlich, neues Dataset versioniert (v1, v2, …), Baseline aktualisiert.
\item \textbf{Observability vom Gateway aus.} Ein Trace: Gateway Request → Provider Call → Cache Hit/Miss → Fallback? → Response. Alles in einem Trace (Observability-Stack: Kapitel 15).
\end{enumerate}

% ============================================================
% ANTI-PATTERNS
% ============================================================
\section{Anti-Patterns}

\begin{itemize}
    \item \textbf{Single-Provider Hardcoding.} \texttt{openai.ChatCompletion.create()} direkt im Business Code. Fix: Gateway Abstraction Layer.
    \item \textbf{In-Memory Circuit Breaker.} State im Prozess → bei Restart/Scale-out verloren. Fix: Redis-backed, distributed State.
    \item \textbf{Kein Cost Budget pro User/Tag.} Ein Bug im Agent-Loop → \$50k/Tag. Fix: Token Bucket Rate Limiter (User + Global) im Gateway.
    \item \textbf{Prompt Changes = Code Deploy.} Prompt im Python String → Jede Änderung = PR + Deploy. Fix: Prompt Templates in DB/Config + Feature Flag.
    \item \textbf{Canary ohne Quality Gate.} 5\% Traffic auf neuen Prompt, aber kein Eval vorab. Fix: Pre-Flight Eval gegen Golden Dataset.
    \item \textbf{Auto-Rollback nur auf Error Rate.} Latenz, Cost, Refusal Rate, Quality Score werden ignoriert. Fix: Multi-Metric Rollback Config.
    \item \textbf{Docs manuell pflegen.} "Wir aktualisieren Docs beim nächsten Sprint." Nie passiert. Fix: Self-Healing Pipeline.
    \item \textbf{Golden Dataset nie aktualisieren.} Eval läuft gegen 6 Monate alte Daten. Fix: Weekly Auto-Eval Pipeline aus Prod Logs.
    \item \textbf{Kein Request Hedging bei Tail Latency.} p99 = 30s, User wartet. Fix: Hedging nach p99-Threshold an Backup-Provider.
    \item \textbf{Alle Traffic auf ein Modell.} Kein Tiering, kein Cost-Aware Routing. Fix: Routing Logic mit Tags (fast/cheap/premium/reasoning).
\end{itemize}

% ============================================================
% ZUSAMMENFASSUNG
% ============================================================
\begin{zusammenfassungEnv}
\begin{itemize}[leftmargin=*]
    \item \textbf{LLM Gateway} ist die zentrale Abstraktionsschicht: Unified API, Resilience (Circuit Breaker, Fallbacks, Rate Limits), Observability, Cost Control.
    \item \textbf{AI Feature Flags} entkoppeln Prompt/Model/Routing-Änderungen von Code-Deployments. Quality Gates (Eval-Score) + Auto-Rollback (Multi-Metric) machen Rollouts sicher.
    \item \textbf{Self-Healing Documentation} schließt die Lücke zwischen Production-Verhalten und Nutzer-Dokumentation. Embedding-Drift-Detection + LLM-Generation + PR-Automation.
    \item \textbf{Auto-Eval Dataset Builder} verwandelt Production Logs in aktuelle Golden Datasets. HDBSCAN Clustering + LLM-as-Judge Labeling ersetzt manuelles Labeling.
    \item \textbf{Architektur-Prinzip:} External Dependencies (Provider) sind unzuverlässig → Gateway absorbiert Volatilität. Probabilistische Outputs → Continuous Evaluation. Kosten pro Request → Cost-Aware Routing + Budgets.
\end{itemize}
\end{zusammenfassungEnv}


\par\mbox{}\par
\begin{merkeEnv}
\begin{enumerate}
    \item \textbf{Gateway = Infrastruktur, nicht Feature.} Baue es zuerst. Ohne Gateway hast du Vendor Lock-in, keine Resilience, keine Cost Control.
    \item \textbf{Feature Flags für AI $\neq$ Feature Flags für Code.} Sie steuern probabilistisches Verhalten. Brauchen Quality Gates (Eval) und Auto-Rollback (Metrics).
    \item \textbf{Production Logs > Synthetische Eval-Daten.} Echte User-Queries decken Edge Cases ab, die du dir nicht ausdenkst. Automatisiere die Pipeline.
    \item \textbf{Self-Healing Docs zahlen sich aus.} Drift Detection + LLM Update + Auto-PR hält Docs aktuell ohne manuellen Aufwand. Threshold 0.85 Cosine Similarity.
    \item \textbf{Cost-Aware Routing ist der Hebel.} 70~\% Traffic auf günstiges Modell, 30~\% auf Premium für Quality — spart 60–80~\% Kosten bei gleicher Qualität.
    \item \textbf{Auto-Rollback auf 4+ Metriken.} Error Rate, P99 Latency, Cost/Query, Refusal Rate. Eine Metrik alleine blendet Regressionen aus.
\end{enumerate}
\end{merkeEnv}

% ============================================================
% PRAXISPROJEKT
% ============================================================
\section{Praxisprojekt — Produktionsreifer LLM Gateway mit Feature Flags \& Self-Healing Docs}

\textbf{Ziel:} Ein vollständiger LLM Gateway Service mit Multi-Provider Routing, AI Feature Flags (Canary + Quality Gate + Auto-Rollback), und Self-Healing Documentation Pipeline — deployed via Docker Compose / Kubernetes.

\subsection*{Aufgaben}

\begin{enumerate}[leftmargin=*]
    \item \textbf{Gateway Core:} Implementiere \texttt{LLMGateway} Klasse (LiteLLM + Redis Circuit Breaker + Fallback Chains + Token Bucket Rate Limiter + Cost Tracking). Endpoint: \texttt{POST /v1/chat/completions} (OpenAI-kompatibel).
    \item \textbf{Model Registry:} Konfiguriere 4 Modelle: Premium-Modell (z.B. GPT-4o), Fast/Cheap-Modell (z.B. GPT-4o-mini), Reasoning-Modell (z.B. Claude Sonnet 4), Self-hosted 70B-Modell (z.B. Llama-3.3-70b). Tags: \texttt{fast}, \texttt{cheap}, \texttt{reasoning}, \texttt{self-hosted}.
    \item \textbf{Fallback Chains:} Premium → Reasoning → Fast/Cheap → Self-hosted 70B. Circuit Breaker Threshold: 5 Failures, Recovery: 30s.
    \item \textbf{Rate Limiting:} User: 60 req/min, 100k tokens/hour. Global: 1000 req/min. Redis-backed Token Bucket.
    \item \textbf{AI Feature Flags:} Integriere Unleash. Definiere 3 Flags: \texttt{support-bot-prompt-v2}, \texttt{model-premium-to-reasoning}, \texttt{enable-semantic-cache}. Je Flag: Eval Dataset, Min Score Threshold, Rollout Steps [5, 25, 50, 100], Auto-Rollback Metriken.
    \item \textbf{Quality Gate:} Pre-Flight Eval gegen Golden Dataset (Langfuse/JSONL). Metric: Faithfulness/Relevance. Threshold: 90~\% of Baseline.
    \item \textbf{Auto-Rollback Monitor:} Hintergrund-Task prüft alle 60s Prometheus Metriken. Bei Violation → Flag Percentage auf 0, Alert an Slack/PagerDuty.
    \item \textbf{Self-Healing Docs Pipeline:} GitHub Action (Daily 03:00 UTC). Observability Traces → Embedding Similarity (text-embedding-3-large) vs. Docs → Threshold 0.85 → LLM Doc Update → Vale + Markdownlint + LLM-Judge → PR → Auto-Merge on Green.
    \item \textbf{Auto-Eval Dataset Builder:} Weekly Job. Observability Logs (7d) → HDBSCAN Clustering → Medoid Sampling (200 Samples) → LLM-as-Judge (GPT-4o, n=3) → Quality Filter (Confidence > 0.8, Consistency 3/3) → Export JSONL v\{N+1\}.
    \item \textbf{Observability:} Prometheus Metriken (Requests, Latency, Cost, Errors, Fallback Rate, CB State). Grafana Dashboard: Overview + Per-Model + Per-Flag. Observability-Stack (Kapitel 15) für jeden Request.
    \item \textbf{Deployment:} Docker Compose (Gateway, Redis, Prometheus, Grafana, Unleash, Observability-Stack). Kubernetes Manifests (Deployment, Service, HPA, ConfigMap, Secret).
\end{enumerate}

\subsection*{Erfolgskriterien (Definition of Done)}

\begin{itemize}
    \item \checkmark Gateway returned OpenAI-kompatible Responses für alle 4 Provider.
    \item \checkmark Circuit Breaker öffnet nach 5 Fehlern, Fallback triggert, Half-Open nach 30s, Closed nach 3 Successes.
    \item \checkmark Rate Limiting blockiert User bei 61. Request/Min mit 429.
    \item \checkmark Feature Flag \texttt{support-bot-prompt-v2} rollt 5~\% → 25~\% → 50~\% → 100~\% mit Quality Gate (Eval Score $\ge$ 0.9 × Baseline).
    \item \checkmark Auto-Rollback triggert bei künstlich erzeugter Error Rate > 5~\% (Chaos Test: inject 10~\% Errors).
    \item \checkmark Self-Healing Pipeline erkennt Drift (geänderte FAQ Antwort), erstellt PR, merged nach Checks.
    \item \checkmark Auto-Eval Builder erzeugt v2 Dataset mit 200 Samples, Quality > 0.85, Export nach JSONL.
    \item \checkmark Grafana Dashboard zeigt Live: Cost/Query, P50/P99 Latency, Error Rate, Fallback \%, CB States.
    \item \checkmark Kubernetes Deployment: HPA (CPU > 70~\%), Readiness Probe (Gateway Health + Provider Reachability), Rolling Update (maxSurge=1, maxUnavailable=0).
\end{itemize}

\subsection*{Erweiterung (Bonus)}

\begin{itemize}
    \item Request Hedging: Nach p99 Latency (3s) Duplicate Request an Backup-Provider.
    \item Semantic Routing: Embedding Classifier (MiniLM) für Intent Detection (<10ms).
    \item Cost-Aware Routing: User Budget \$5/Tag → Gateway routet automatisch zu günstigerem Modell bei Budget-Erschöpfung.
    \item Shadow Traffic: 100~\% Traffic an Current + Candidate, Metrik-Vergleich ohne User Impact.
    \item Multi-LoRA Hot-Swap via Gateway (Vorbereitung Ch18).
\end{itemize}

% ============================================================
% INTERVIEWFRAGEN
% ============================================================
\section{Interviewfragen}

\subsection*{F1: Warum reicht ein In-Memory Circuit Breaker nicht für LLM Gateways in Produktion?}
\textbf{Antwort:} In-Memory State geht bei Restart/Scale-out verloren. Bei 3 Replicas hat jeder eigenen State → Inconsistente CB-Entscheidungen. Verteilter CB (Redis) teilt State: Alle Instanzen sehen gleichen State, Recovery koordiniert, keine "Split Brain" Failures.

\subsection*{F2: Was ist der Unterschied zwischen klassischem Feature Flag und AI Feature Flag?}
\textbf{Antwort:} Klassische Flags steuern deterministischen Code-Pfad (An/Aus). AI Flags steuern probabilistisches Verhalten (Prompt, Model, Temperature, Retrieval Config). Deshalb brauchen AI Flags: (1) Quality Gates (Eval vor Rollout), (2) Auto-Rollback auf Business Metriken (Cost, Latency, Quality), (3) Canary/Shadow Traffic Support, (4) Gradual Rollout mit Metrik-Monitoring pro Step.

\subsection*{F3: Wie funktioniert HDBSCAN und warum ist es besser als K-Means für Production Log Clustering?}
\textbf{Antwort:} HDBSCAN baut Hierarchie von Dichte-basierten Clustern, extrahiert stabile Clusters über Dichte-Ebenen (Excess of Mass). Vorteile: Kein k nötig, findet variable Dichten (FAQs dicht, Edge Cases sparse), markiert Noise explizit (Outliers = seltene/fehlerhafte Queries), robust gegen Outliers. K-Means zwingt kugelförmige Clusters, fixed k, sensitiv zu Outliers.

\subsection*{F4: Was sind die 4 kritischen Metriken für Auto-Rollback bei AI Feature Flags?}
\textbf{Antwort:} (1) \textbf{Error Rate} > 5~\% (5min Window) — harte Fehler. (2) \textbf{P99 Latency} > 3× Baseline (10min Window) — Tail Degradation. (3) \textbf{Cost/Query} > 1.5× Baseline (60min Window) — Kostenexplosion. (4) \textbf{Refusal Rate} > 10~\% (10min Window) — Modell verweigert (Safety/Alignment Drift). Optional: Quality Score (Faithfulness/Relevance) < 0.8 × Baseline.

\subsection*{F5: Wie verhindert man, dass Self-Healing Docs Halluzinationen in die Dokumentation schreiben?}
\textbf{Antwort:} Mehrstufige Validation: (1) \textbf{Source Grounding:} LLM Prompt verbietet Halluzinationen — "Nur Fakten aus QA-Pairs übernehmen". (2) \textbf{Vale/Markdownlint:} Style, Links, Struktur. (3) \textbf{LLM-as-Judge Factuality Check:} Zweiter LLM Call prüft: "Stimmt generierter Content mit QA-Pairs überein?" (4) \textbf{Human-in-the-Loop:} PR requires Reviewer Approval (Docs-Team) vor Auto-Merge. Auto-Merge nur bei allen Checks grün.

\subsection*{F6: Was ist Request Hedging und wann lohnt es sich?}
\textbf{Antwort:} Nach p99-Latency-Threshold (z.B. 3s) wird identischer Request an Backup-Provider gesendet. Erste Response gewinnt. Kostet ~1–5~\% Extra API Calls (nur Tail Requests). Lohnt sich wenn p99 > 10s und User-facing Latency kritisch ist (Chat, Search). Nicht lohnen bei Batch/Async Workloads.

\subsection*{F7: Wie designst du Cost-Aware Routing ohne Quality-Verlust?}
\textbf{Antwort:} (1) \textbf{Tag Models:} cheap/fast/premium/reasoning/self-hosted. (2) \textbf{Task Classification:} Simple QA → cheap; Complex Reasoning → premium; Code → code-specialized. (3) \textbf{Budget Enforcement:} User/Day Budget im Gateway → bei Erschöpfung Auto-Downgrade zu cheap Model. (4) \textbf{Quality Sampling:} 10–30~\% Traffic immer auf Premium für Quality Monitoring. (5) \textbf{Continuous Eval:} Täglicher Auto-Eval Report vergleicht Cheap vs. Premium Quality Gap.

\subsection*{F8: Warum ist Shadow Traffic (Mirroring) sicherer als direkter Canary?}
\textbf{Antwort:} Shadow Traffic sendet 100\% Requests an Current + Candidate parallel. User sieht nur Current Response. Candidate Metriken (Latency, Cost, Quality, Errors) werden gesammelt ohne User Impact. Erst nach 24h Shadow-Analyse (keine Regressionen) startet Canary Rollout. Canary ohne Shadow risikiert echte User bei Regressionen.

% ============================================================
% WEITERFÜHRENDE RESSOURCEN
% ============================================================
\section{Ressourcen}

\subsection*{Gateway \& Multi-Provider Routing}
\begin{itemize}[leftmargin=*]
    \item \textbf{LiteLLM Documentation:} \url{https://docs.litellm.ai} — Gateway Implementation Reference, 100+ Provider, Fallbacks, Budgets, Callbacks.
    \item \textbf{Portkey Blog — "Building a Production LLM Gateway":} \url{https://portkey.ai/blog} — Architecture Patterns, Caching, Fallbacks, Virtual Keys.
    \item \textbf{Helicone Docs:} \url{https://docs.helicone.ai} — Observability-first Gateway, Prompt Management, Caching.
\end{itemize}

\subsection*{Feature Flags \& Progressive Delivery}
\begin{itemize}[leftmargin=*]
    \item \textbf{Unleash Docs — "Feature Flags for AI Applications":} \url{https://docs.getunleash.io} — Strategies, Gradual Rollout, Kill Switches, AI-specific Use Cases.
    \item \textbf{Argo Rollouts / Flagger:} Kubernetes-native Progressive Delivery — AnalysisTemplates, Metric Providers (Prometheus, Datadog), Automated Rollback.
\end{itemize}

\subsection*{Evaluation, Clustering \& Observability}
\begin{itemize}[leftmargin=*]
    \item \textbf{HDBSCAN Paper:} Campello et al. (2013) "Hierarchical Density Estimates for Data Clustering" — Theoretische Grundlage für Variable-Density Clustering.
    \item \textbf{LLM-as-Judge:} Zheng et al. (2024) "Judging LLM-as-a-Judge with MT-Bench and Chatbot Arena" — GPT-4 Judge korreliert 0.85+ mit Human Preference; Konsistenz durch n=3 + Confidence.
    \item \textbf{Google — "Continuous Evaluation of Language Models in Production" (2024):} Shadow Evaluation + Canary + Auto-Rollback = Production Quality ohne Manual Gates.
    \item \textbf{Anthropic — "Prompt Caching \& Evaluations":} \url{https://docs.anthropic.com} — Prompt Caching (90~\% Prefill Latency Reduction), Eval Framework für Regression Testing.
    \item \textbf{Databricks — "Building LLM Applications with MLflow \& Gateway" (2024):} MLflow Gateway für Multi-Provider Routing + Evaluation Tracking + Model Registry Integration.
\end{itemize}

\subsection*{Self-Healing Documentation \& CI/CD}
\begin{itemize}[leftmargin=*]
    \item \textbf{GitHub Actions — "Automated Documentation Updates" (2024):} Actions + LLM für Doc-Generation aus Code/Logs — PR-basiert, review-gated, auditierbar.
    \item \textbf{Vale + Markdownlint:} Style-Guide Enforcement + Syntax-Check für generierte Docs.
\end{itemize}

\subsection*{Tool-Vergleichstabellen}
\begin{table}[H]
\centering
\begin{tabularx}{\linewidth}{lXXX}
\toprule
\textbf{Tool} & \textbf{Typ} & \textbf{Stärken} & \textbf{Schwächen} \\
\midrule
LiteLLM & Python Proxy / SDK & 100+ Provider, OpenAI-kompatibel, Fallbacks, Rate Limits, Budgets, Callbacks & Single-Process (skaliert via Replicas), keine native K8s CRDs \\
Portkey & Managed Gateway (SaaS/Self-hosted) & Analytics, Caching, Fallbacks, Virtual Keys, Prompt Templates, Canary & Closed Source Core, Pricing bei Scale \\
Helicone & Observability + Gateway & Open Source, Logging, Caching, Rate Limits, Prompt Management & Weniger Routing-Features als LiteLLM \\
Langfuse Gateway & Observability-First & Tracing, Evals, Prompt Management, Datasets & Gateway-Features noch jung \\
\bottomrule
\end{tabularx}
\caption{LLM Gateway / Multi-Provider Router Vergleich}
\label{tab:ch17_gateway_tools}
\end{table}

\begin{table}[H]
\centering
\begin{tabularx}{\linewidth}{lXXX}
\toprule
\textbf{Tool} & \textbf{Typ} & \textbf{Stärken} & \textbf{Schwächen} \\
\midrule
Unleash & Open Source (Self-hosted) & Enterprise Features kostenlos, Strategies (UserID, IP, Custom), API, Webhooks & UI weniger poliert als LaunchDarkly \\
LaunchDarkly & SaaS & Best-in-class UI, Experimentation, Data Export, Integrations & Teuer (\$\$\$), Vendor Lock-in \\
GrowthBook & Open Source / SaaS & A/B Testing Focus, SQL-basiert, Visual Editor & Feature Flag Features Basis \\
Flagsmith & Open Source / SaaS & Simple, Segments, Multivariate Flags & Weniger Enterprise Features \\
\bottomrule
\end{tabularx}
\caption{AI Feature Flag Tools Vergleich}
\label{tab:ch17_flag_tools}
\end{table}

\begin{table}[H]
\centering
\begin{tabularx}{\linewidth}{lXXX}
\toprule
\textbf{Komponente} & \textbf{Zweck} & **Empfohlen** & \textbf{Alternative} \\
\midrule
Log Source & Production Traces & Langfuse / Helicone & LangSmith, Custom \\
Clustering & Topic Discovery & HDBSCAN (cosine) & K-Means (fixed k) \\
Sampling & Representative Selection & Medoid (robust) & Centroid (schneller) \\
LLM Judge & Auto-Labeling & GPT-4o + Claude + Gemini (n=3) & Single Judge (Bias) \\
Validation & Quality Gates & Confidence > 0.8 + 3/3 Consistency & Threshold only \\
Export & CI/CD Integration & JSONL (OpenAI Eval) / HF Dataset & CSV \\
\bottomrule
\end{tabularx}
\caption{Auto-Eval Dataset Builder Pipeline Komponenten}
\label{tab:ch17_autoeval_components}
\end{table}

% ============================================================
% NÄCHSTER SCHRITT
% ============================================================

\par
\begin{autornotizEnv}
Nächstes Kapitel: Production Checklist — Der Pre-Flight-Check für jeden Deploy. Design, Development, Security, Production, Evaluation, Governance, Kosten — alles auf einem Blatt. Dein Checklist-Workflow für Zero-Surprise-Deployments.
\end{autornotizEnv}

\textbf{Nächster Schritt:} Mit Gateway, Feature Flags, Self-Healing Docs und Auto-Eval hast du die Produktions-Infrastruktur für LLM-Systeme beherrscht. Der letzte Schritt ist die Production Checklist: Dein Pre-Flight-Check für jeden Deploy.

\endinput